\subsection{Measurement, contextuality, and effective descriptions}
  \label{subsec:contextuality}

  Recent work in quantum foundations has emphasized that nonlocal and contextual
  correlations may be understood as obstructions to constructing globally consistent
  probabilistic models from local measurement contexts~\cite{AbramskyBrandenburger2011}.
  Within ontological-model frameworks, observable statistics are likewise interpreted
  as coarse-grained descriptions of an underlying ontic state space~\cite{Spekkens2007}.
  These approaches identify contextuality as a fundamental feature of quantum theory,
  reflecting the impossibility of assigning outcome values independently of the
  measurement context.

  The mechanism discussed here is complementary in spirit, but identifies a more
  primitive structural origin for the failure of Bell factorizability.
  Rather than attributing contextuality to context-dependent properties of observables,
  we trace it to the non-injective mapping between underlying configurations and
  observable descriptions.
  When the effective description has a finite capacity to encode distinctions present
  at the underlying level, distinct configurations are necessarily identified, and no
  globally consistent assignment of independent local outcomes is possible.

  From this perspective, contextuality appears as a manifestation of descriptive
  non-injectivity.
  In particular, the impossibility of assigning non-contextual value assignments
  can be traced to the fact that observable variables do not separate equivalence classes
  of underlying configurations.
  The impossibility of constructing non-contextual models reflects the fact that
  observable variables do not provide a complete specification of the underlying
  configuration.
  Bell-type correlations and contextuality therefore share a common structural origin
  in the finite-capacity nature of effective descriptions.

  In regimes where the effective descriptive capacity is large compared to the
  explored configuration space, the mapping becomes approximately injective.
  In this limit, contextuality effects are suppressed, and observable outcomes may be
  treated as effectively independent.
  Classical statistical behavior thus emerges as a regime of approximate injectivity,
  rather than as a fundamentally distinct domain.

  The present work does not propose a modification of quantum mechanics or a new
  ontological framework.
  Its contribution is to isolate a minimal and sufficient structural condition for the
  failure of Bell inequalities and the emergence of contextual correlations.
  Non-injective projection, understood as a finite-capacity effect, provides such a
  condition while preserving both operational locality and empirical adequacy.

  We conclude that Bell inequality violations and contextuality need not be interpreted
  as evidence for fundamental nonlocality, but may instead reflect intrinsic limitations
  of finite-capacity effective descriptions in which observable outcomes are treated as
  independently specifiable local properties.
